\chapter{项目特色与实践价值}
\section{围绕同一订单的四端协作}
项目特色首先体现在业务链路而非页面数量。顾客负责选餐和确认，商家负责接单备餐，骑手负责配送，管理员处理审核与平台管理。四端共享订单事实但拥有不同操作权限，状态推进必须与角色权限一致。这样的结构使课程实践同时涉及前端交互、接口契约、事务处理、状态机与多角色测试。

设计上的价值是让业务分工可见：顾客看到的“等待”应对应商家或骑手可执行的动作，而不是各端独立维护一套状态。实践上的难点则是跨端联调和异常同步，不能用单个页面可点击证明链路已闭合。
\section{外送与到店自取双履约方式}
订单支持外送与到店自取，顾客在结算前选择方式，系统据此处理地址要求、配送费和后续履约步骤。自取不应套用必须填写配送地址或必须创建骑手任务的规则。两种方式让项目从单一路径扩展到条件分支，也使需求、报价、页面提示和订单状态必须同步设计。
\section{AI 点餐与真实业务数据结合}
AI 入口不应只是与外部模型聊天。项目通过意图分流和业务查询，将预算、口味、菜单及本人订单与回答关联。结构化候选商品需要经过商品存在性、价格、库存和归属校验，加入购物车或提交订单仍由确定性的业务接口处理。模型负责理解和表达，服务端负责业务事实与交易规则。

图片和语音为交互提供了额外入口，但能力依赖外部服务及配置。特色展示时应说明可用条件，不将所有入口都展示为无条件可用，也不将智能化等同于放松订单确认和授权要求。
\section{红包积分、看板与地图的业务连接}
红包积分不是孤立的余额显示，需要与报价、支付、退款或取消回退协调；看板需要统一订单量、收入与状态的统计口径；高德地图选择地址后还需保存坐标，并与配送距离及页面展示衔接。三类扩展分别连接资金规则、经营观察与真实地理位置，体现了从基础 CRUD 向业务约束的扩展。
\begin{longtable}{L{.21\linewidth}L{.34\linewidth}L{.36\linewidth}}
\caption{特色模块的价值与验证重点}\\\toprule 模块&用户价值&验证重点\\\midrule\endfirsthead\toprule 模块&用户价值&验证重点\\\midrule\endhead
经营看板&集中了解经营与平台情况&订单状态口径、金额口径、筛选与空数据展示，避免把演示数值当作真实经营成果。\\
红包与积分&让优惠和资产可在结算中使用&服务端重算、扣减顺序、重复支付防护、取消或失败后的资产一致性。\\
高德地图&辅助选择和理解配送位置&坐标与地址保存一致、地图失败回退、定位授权、距离的来源和展示。\\
AI 辅助点餐&降低筛选和表达需求的成本&业务数据依据、权限、调用次数、总耗时、结构化结果验证与用户确认。\\\bottomrule
\end{longtable}
\section{特色的适用边界}
上述内容是本项目的组合特点和课程实践重点，不表示行业首创，也不代表已经达到生产系统规模。订单支付仍使用演示方案，真实并发性能和外部 AI 能力需要专门验证。项目价值在于成员能够解释功能之间的关系、明确失败场景，并在老师反馈和交叉测试中持续修正实现与文档。
